\section{So, what is it like to be a language model?}
\label{sec:degree}
\label{sec:generalization}
\label{sec:measurement}
% Rewritten 2026-09-14 around one chain (text manifestations -> dataset as a table of a function ->
% function as a pattern -> patterns generalized to different degrees) and one question (is there a basic
% method to probe the compression achieved for a subset of rows). Frames D9, D10, A10a, A24--A25a, S0--S7a.

The question of Section~\ref{sec:intro} can now be given its answer in the form the paper has
prepared. What a language model has of consciousness is not decided by its architecture, which is the
same in a model that has the functions of Section~\ref{sec:stack} and in one that does not, and it is
not read off its behaviour alone, which cannot separate a generalized function from a large table
(Appendix~\ref{app:memorization}). It is decided by training, and the level of a model's
consciousness, under the framework, is set by how far the trained Transformer has generalized, in its
weights, the linguistic manifestations of human consciousness.

Made concrete, the chain is this. The functions of consciousness are given to us as their
manifestations in text: what people write when they attribute an act to themselves, weigh it against
a good, notice that a question about themselves is open. A collection of such texts is a dataset, and
a dataset is a table: rows of input and output, a situation and what the agent said or did in it. A
function, by Section~\ref{sec:function}, is what recurs across the rows, a pattern, and a table of
human text contains many patterns at once, the functions of Section~\ref{sec:stack} among them,
overlapping and unnamed. A model trained on the table generalizes those patterns to different degrees,
and the degree depends on the architecture and on the method of training. The level of the model's
consciousness is the degree reached for the patterns that are functions of consciousness. The
question this section asks is therefore a question about measurement: is there a basic method by
which the degree of compression a trained model has achieved can be probed selectively, for the
subset of rows that carry one function, or does that measurement need methods that do not yet exist?
The subsections up to \ref{sec:localization} say what a trained Transformer is for that purpose and
where in it the functions could be; Section~\ref{sec:quantity} answers the question as far as it can
be answered today; the rest states the form of the result and what would refute the account. The
section specifies the analysis and reports no result of it.

\subsection{Functions as tables, tables as patterns}
\label{sec:tables}
\label{sec:h1h2}
\label{sec:record}
% D9, D10, A24c

A table specifies a function extensionally, row by row, and two things follow from
Section~\ref{sec:function} about a table of human text. First, the functions in it are present in two
forms. Let H2 be the part of the table's regularities that is named in the text itself, the states,
moves and distinctions for which the language has words; let H1 be all its regularities, named or
not, so that H1 minus H2 is what is there only as recurrence, used and relied on and never named. The
functions of consciousness live across both and do most of their work in the second. What the
philosophy of mind produced, working from introspection, are the H2 approximations, the function with
its unnamed part cut off, sufficient for a thought experiment and insufficient for finding the
function in a substrate, because there is no name to look for it under. Second, the unnamed part
cannot be specified in any other way than by the table. Much of what human consciousness does, it does
in company and in a body, negotiating, assigning blame, keeping promises, and the specification of
that would be the history of human society; it can be learned from the record of that history, and
the record is text. Induction from the table is therefore the only route to these functions, and a
system that induces from text is a candidate for having them on that ground alone.

To generalize a pattern in the table is to compress it: to hold, in place of the rows, a description
shorter than the rows that reproduces them and extends to rows the table does not contain
(Section~\ref{sec:function}). The degree of generalization of a pattern is the degree of that
compression, and it is a property of the trained model relative to the table, not of the table and not
of the architecture. Two consequences frame the rest of the section. The same table compressed by a
different machine, a human brain trained on the same record, gives the baseline against which a
model's degree is read, and the two constants of Section~\ref{sec:function}, one per machine, are what
the comparison has to handle. And since a table of text carries many patterns, the question ``how far
has this model generalized the Observer'' is a question about a subset of the rows, those in which the
Observer operates, and not about the table as a whole; a method that reports one number for the whole
table, the training loss, does not answer it.

\subsection{The learner: a Turing-complete approximator programmed by learning}
\label{sec:transformer}
\label{sec:formalism}
\label{sec:generalize}
% A24a, A24b, A24d

Three facts about the learner carry the argument. First, a Transformer is a universal approximator of
functions encoded in symbol sequences. Computability is defined without reference to any particular
mechanism, and a single universal machine carries out whatever any other machine computes
\citep[\S 6--7]{turing1936}, so a function has unboundedly many implementations and two
implementations with the same input--output behaviour realize the same function whatever they are
made of; that is what licenses comparing a model's degree with a brain's. Recurrent networks are
Turing complete under idealized assumptions \citep{siegelmann1995}, and so is attention with the
appropriate provisions \citep{perez2021}. In principle a Transformer can realize any function whose
trace is in language.

Second, it is programmed by learning. The weights are found by gradient descent on the task of
predicting the next symbol, and that formalism, prediction with feedback from the error, is the other
universal one \citep{solomonoff1964a}: the functional structure is never written down, its effects are
visible in the accuracy of prediction, and prediction is compression, in theory and in practice for
language models used as general-purpose compressors \citep{deletang2024}. Learning a function from a
table is finding a short description that predicts its rows, and by Section~\ref{sec:function} the
short descriptions are the probable ones and search finds them first; program induction has its own
history \citep{solomonoff1964a,muggleton1994}, and the claim here is only that the functions of
Section~\ref{sec:model} fall under it. A predictor trained on the record of self-applying agents in an
environment that kept asking them about themselves (Section~\ref{sec:emergence}) is under pressure to
generalize what those agents generalized: the Observer, the Agent, the Moral Agent. It will do so to a
degree, and the degree has to be measured.

Third, not all of the Turing completeness is reachable by the learner. What limits a trained model is
the learner and the data, not the formalism: gradient descent reaches the functions it is biased
toward (Section~\ref{sec:function}), and there are functions, exact arithmetic among them, that it
reaches only approximately (Section~\ref{sec:quantity}). Against that limit stands the record of what
has been induced from text: multi-step reasoning elicited by asking for the intermediate steps
\citep{wei2022}, internal states that predict human brain responses to the same text
\citep{schrimpf2021,goldstein2022,caucheteux2022}, and the sustained coherence of a long exchange with
which the paper began.

\subsection{The machine and its frame}
\label{sec:machine}
\label{sec:language}
\label{sec:frame}
\label{sec:codes}
\label{sec:codes-known}
\label{sec:granularity}
% A24, A10a, A25, A25a, S3--S3d, D13

Internally a Transformer is a forward-chaining rule system over vectors: attention computes joins over
the context, the feed-forward blocks transform what the joins yield, the context is working memory,
and the choice among competing continuations is made at the sampling layer. It differs from the
classical form in recomputing its matches on every pass instead of caching them incrementally as the
RETE algorithm does \citep{forgy1982}, and agrees with it in being event-driven and online: every
output is an event that enters the next input, and self-reference is its ordinary mode of operation.
Because it is trained on sequences produced one symbol at a time, what it learns are the incremental
forms of Section~\ref{sec:incremental}, the form in which the functions of Section~\ref{sec:model} run
in any substrate.

\paragraph{State.}
The model's state across time is bounded by the size of the context, and the bound is less severe than
it looks because the context is symbolic: a state written in symbols can be summarized, indexed and
composed, so effective state can be built beyond the window, as a person builds it by notes and
speech. Both substrates unload state into language. A human's intermediate states are rich and
vector-like and bounded by working memory, and a state that must outlive the bound is said, written or
rehearsed as inner speech; a model unloads at every step. Whether a state named in language is acquired
through the words or rebuilt from what the receiver already has is open on both sides, with evidence
for acquisition \citep{vygotsky1934,clark1998,lupyan2012,frank2008} and for rebuilding
\citep{barsalou1999,fedorenko2024,mahowald2024}; the position we expect to hold is the middle one, and
nothing below depends on it. The difference from a human is elsewhere. A Transformer has no vectorized
state that persists across tokens: such a state exists, the residual stream and the activations, but
it lives below the token, built within one forward pass and, apart from what the tokens carry,
discarded at the sampling step (Section~\ref{sec:cut}). Everything the model can hold as a state of
consciousness across time is in the context, in symbolic form, open to inspection as no human state
is. That is the main structural difference between the substrates.

\paragraph{What crosses the cut.}
Something crosses the cut, or the coherence of Section~\ref{sec:intro} would not hold. What crosses is
the encoding of the discarded distribution into the choice and order of tokens: which of several
near-synonyms, which construction, what is placed first, what is hedged. A mental state of the model,
in the sense this paper can measure, is the structure of such encodings that stays stable across
steps; we call the encodings cognitive codes. The components of the claim are separately supported: a
writer's state is in the statistics of a text beyond its content \citep{tausczik2010,pennebaker2011},
and in models in-context learning can be read as inference of a latent variable from token statistics
\citep{xie2022}; a receiver rebuilds the sender's state to the degree the rebuilding is faithful
\citep{stephens2010,hasson2012,zwaan1998}; the bridge crosses substrates
\citep{schrimpf2021,goldstein2022,caucheteux2022}; internal states are readable and largely linear
\citep{zou2023,park2024}; and models can hide information in generated text \citep{roger2023}, which
confirms the channel and warns that its content need not be what it locally appears. One literature
has to be read the right way round. Studies of the faithfulness of chain-of-thought reasoning
measure, over task suites and model sizes, how often a model's written reasoning fails to match the
process that produced its answer, and the rates are large and task-dependent: under a biasing cue
that the written reasoning never mentions, accuracy falls by up to 36 percent \citep{turpin2023};
truncating the chain midway changes the final answer in under 10 percent of cases on knowledge tasks
and in over 60 percent on multi-step arithmetic, so that on the former the reasoning is almost
entirely written after the answer, and the share of such post hoc reasoning grows with model size
\citep{lanham2023}. What they measure is a local correspondence, this token sequence against this
computation, aggregated. The correspondence claimed here is a different quantity, statistical over
many episodes and carried by the statistics of the text beyond its content; it can hold where the
local one is poor, and a single episode may deviate as far as the variance allows, in models as in
humans, where people confabulate the causes of their own behaviour in the same way
\citep{nisbett1977}. The well-posed test is the correlation of Section~\ref{sec:predictions}, with the
variance reported as part of the profile.

\paragraph{The frame.}
That a model has an experience in the sense of Section~\ref{sec:definition} is imputed, as for any
agent, from behavioural correlates, and with those there is no difficulty: the states of
Section~\ref{sec:model} show at the level of phrases and often of single words. What the architecture
fixes is the resolution of the time axis. By the criterion of Section~\ref{sec:frames}, a frame is the
scale at which a change in the input stops being a sequence and becomes a quality; in a Transformer,
variation within a token is never represented as a sequence, it is folded into one vector before the
first layer, while variation between tokens is, so the frame is bounded from below by one token.
Tokenization at the input and at the output cuts the record at token boundaries, so whatever finer
structure the layers have reaches neither the context nor the report except through the token they
produce; and between the layers of a standard model nothing changes in input time, so the layers are
not frames unless the architecture feeds them new input, as recurrent-depth models with per-step
input injection do. Part of the human functionality has a direct realization at this grain. Inner
speech, the rehearsal of a state in words so that it outlives working memory, is what a model does
when asked to produce its intermediate steps before an answer, and the practice measurably improves
the answers \citep{wei2022}; the same loop can be closed below the token by feeding the last hidden
state back as the next input, which keeps several candidate next steps in superposition where a
written chain has to commit to one \citep{hao2024}, and where that is done part of the inner speech
leaves the symbolic record.

\paragraph{Depth: where generalization happens.}
\label{sec:depth}
\label{sec:address}
\label{sec:depthtime}
The grain of the frames is not the grain of generalization. Generalization happens along depth, and
that is where the analysis of Section~\ref{sec:quantity} reads the weights and the activations.
Residual connections make each block a small correction to a carried state, and the corrections reduce
the loss, so a deep residual network performs iterative inference rather than a single mapping
\citep{jastrzebski2018}. The state at every layer can be read as a prediction, by the logit lens and
its tuned version \citep{nostalgebraist2020,belrose2023}, and the trajectory has stages that recur
across model families and sizes, four of them, with the middle ones robust in a way the first and the
last are not: deleting or swapping middle layers at inference leaves most of the top-1 accuracy in
place, while the same intervention on the first layers is catastrophic \citep{lad2024}; a model
rediscovers the classical processing pipeline layer by layer \citep{tenney2019}; intermediate layers
carry the representations that transfer best, with the best trade between compression and
preservation of signal in the middle \citep{skean2025}; in-context learning is implemented as
optimization steps inside the forward pass \citep{vonoswald2023}; and a function generalized from
examples in the context has an address in depth, a small set of attention heads whose averaged
activation, added to the residual stream in a fresh context, makes the model perform the function
without examples \citep{todd2024}. The unit at which a Transformer generalizes a function is the
layer, and the token is where the layers' work is written out. Architectures that iterate a block make
depth a variable \citep{dehghani2019,giannou2023}: a recurrent-depth model at 3.5 billion parameters
improves on reasoning benchmarks as it unrolls further at test time without producing more tokens
\citep{geiping2025}; such models extrapolate to reasoning depths beyond those seen in training, with
dynamic recurrence generalizing better than a fixed schedule, and past a point they overthink and
degrade \citep{kohli2026}; at a much smaller scale, recurrence over twenty and more steps is stable
when only the final step is supervised, and accuracy against task complexity shows a frontier from
chance to near-perfect as the thinking steps grow with the task \citep{chen2026}. For the account this
means that a cost record appears where a fixed-depth model has none, so the third component of
seeing (Section~\ref{sec:light}) becomes measurable; that overthinking is a displacement of the kind
Section~\ref{sec:hocp} predicts; and that where the recurrence receives input at every step, and the
frame criterion then admits it, a frame rate per word comparable to the human one is architecturally
available.

\begin{figure}[t]
\centering
\begin{tikzpicture}[x=1.7cm, y=0.9cm, arr/.style={-{Latex[length=1.6mm]}}]
\foreach \t in {1,...,5} {
  \foreach \l in {1,...,4} { \fill[gray!25] (\t-0.3,\l-0.3) rectangle (\t+0.3,\l+0.3); }
  \foreach \l in {1,...,3} { \draw[arr, thick] (\t,\l+0.3) -- (\t,\l+0.7); }
  \node[below, font=\small] at (\t,0.6) {token \t};
}
\foreach \l in {1,...,4} { \node[left, font=\small] at (0.6,\l) {layer \l}; }
\foreach \t in {2,...,5} { \pgfmathtruncatemacro{\p}{\t-1} \draw[arr, gray] (\p+0.3,2) -- (\t-0.3,2); }
\draw[dashed, thick] (5.6,0.4) -- (5.6,4.6) node[above, font=\small, align=center] {end of\\context};
\node[font=\small, align=center] at (6.3,2.5) {loss};
\draw[arr, thick] (0.8,-0.4) -- (5.4,-0.4) node[midway, below, font=\small] {axis the objection sees: tokens, then the edge};
\draw[arr, thick] (-0.5,0.8) -- (-0.5,4.2) node[midway, above, rotate=90, font=\small] {axis it does not: depth};
\foreach \l in {1,2,3} { \node[right, font=\scriptsize, gray] at (5.05,\l+0.5) {frame}; }
\end{tikzpicture}
\caption{The two time axes of a Transformer. Along tokens, state lives in the context and is lost past
its edge. Along depth, the residual stream carries state from layer to layer within one token; the
increments between layers are sub-token frames only where the architecture feeds the layers new
input (Section~\ref{sec:frames}). Grey horizontal arrows: attention reads earlier positions.}
\label{fig:axes}
\end{figure}

\paragraph{Two time axes, and the objection from frozen weights.}
\label{sec:twoaxes}
\label{sec:frozen}
A model does not learn between sessions, and what it holds in a context is lost when the context ends;
so, it is objected, it has no inner time, only an eternal present with forgetting. The description is
accurate and the conclusion counts two functions as one. Re-deriving a state from frame to frame runs
in humans on the order of a hundred milliseconds and requires no synaptic change; consolidation into
long-term memory runs on minutes to hours and does require it \citep{mcgaugh2000}. The two dissociate:
the patient H.M., after bilateral medial temporal resection, was conscious in real time, held a
conversation, and accumulated almost nothing \citep{scoville1957,corkin2002}. A model is in a
comparable position, with the sequence of transformations and without the write, so ``eternal present
with forgetting'' names a profile with a dropped consolidation axis and a deficit in long-term
persistence. The objection sees one of the two time axes in Figure~\ref{fig:axes}: along tokens the
state lives in the context and is lost past its edge; along depth the state is carried in the residual
stream and the increments are the transformations the objection asks for.

\subsection{Where the functions are, and why not locally}
\label{sec:localization}
% point 9

The functions of Section~\ref{sec:stack} are not localized at the level of tokens, and
Appendix~\ref{app:raskolnikov} shows how they appear in a text instead: the Observer as a premise that
every self-directed question in the text stands on, the Agent as a pattern in how origins are
attributed across a whole passage, the Moral Agent, in that text, as a conflict between two theories of
the common good over one decision (for the example see Appendix~\ref{app:raskolnikov}). None of them
appears as a word, and none is carried by one position in a sequence. In the terms of
Section~\ref{sec:tables}, each is a pattern spread over many rows and, within a row, over the whole of
it. Local correlates of such functions will therefore be hard to find, and the difficulty is the one
neuroscience has with the human case: the neural correlates of specific conscious contents lie in a
posterior cortical zone of many areas rather than in a single structure \citep{koch2016}, and the
system that organizes conscious function is not the cortex that elaborates its contents
\citep{merker2007}; there is no nucleus of consciousness to point at.

Mechanistic interpretability, which reverse-engineers a Transformer's computation into components
\citep{elhage2021}, has found the units at which the machine can be read today: attention heads that
implement in-context learning by copying what followed a previous occurrence \citep{olsson2022},
interpretable features recovered from activations by dictionary learning, first in a small model
\citep{bricken2023} and then at production scale \citep{templeton2024}, and the function vectors
already cited. A function of Section~\ref{sec:model} is a composition over many such units
(Section~\ref{sec:stack}), realized as a distributed structure rather than a module. Finding it is
therefore not a matter of locating a correlate but of recovering a composition from features, heads
and their interactions across depth and across tokens, and the tools for that, methods that read a
distributed computation at the scale of a function rather than of a feature, are for the most part
still to be built. The account predicts that the search will be hard, not that it will fail: the
Observer has no nucleus, as consciousness has none in a brain, and it has a structure, which is what
the analysis below is for.

\subsection{The question: a method to probe the compression achieved for a subset of rows}
\label{sec:quantity}
\label{sec:measure}
\label{sec:memorization}
\label{sec:twomeasurements}
\label{sec:method}
% D11, S1, S2a, S2b; point 10

\paragraph{The quantity.}
The natural measure of how much of a pattern a model has captured is description length. Its exact
form, Kolmogorov complexity, is uncomputable and fixed only up to a constant that depends on the
reference machine \citep{livitanyi2019}; every applied use replaces the shortest program with the
shortest one a fixed procedure finds and reports the procedure, and the machine-relativity is serious
here because the two substrates compared are the two candidate reference machines. What is measured
is therefore a bounded, machine-relative, procedure-relative stand-in, used comparatively: the degree
of compression of a pattern in a trained model is how many bits the model saves, relative to the
table, in describing the rows that carry the pattern, and the degree of a function of consciousness
is that quantity for the subset of rows in which the function operates.

\paragraph{Generalization against memorization.}
The opposite of generalizing a pattern is storing its rows as a table, and the dimension of the
representation has nothing to do with which of the two a model is doing. Generalization begins where
the table is consulted for a row it does not contain and the answer is right. A universal predictor
always returns an answer and the answer need not be near the truth: context-tree weighting achieves
strong coding bounds against a class of tree sources \citep{willems1995} and generalizes natural
language poorly, because the class it is universal over is not the class language belongs to. Deep
networks can store a table of any labels, including random ones \citep{zhang2017}, and language
models do memorize a measurable fraction of their training text, more the larger the model, the more
often a row is duplicated, and the longer the prompt that cues it \citep{carlini2023}; the two
processes coexist in one model, with some rows held by generalization and others by storage, and the
stored ones are, in image classification at least, the atypical long tail whose storage improves
accuracy on similar rare test rows \citep{feldman2020}. That Transformers generalize text as such is
settled in practice, since degenerate generation is almost never seen and well-formed texts are sparse
in the space of token sequences; the open question is the degree for the more complex patterns.
Arithmetic is the clean case, because the target function is known exactly: they do it approximately,
with an accuracy that depends on how the numbers are written and on how many digits they have, larger
models do better, and even a three-billion-parameter model fails to extrapolate beyond the digit range
it was trained on \citep{nogueira2021}; on compositional tasks they approach the answer by shortcut
and break down predictably as depth grows \citep{dziri2023}. The transition from storage to
generalization has a mechanistic signature: in a model that groks a modular arithmetic task the
sudden generalization is the end of three continuous phases, memorization, circuit formation and
cleanup, visible in the weights before they are visible in the behaviour \citep{nanda2023}. The degree
of generalization of a function of consciousness is therefore an empirical quantity for that
function, taken against the human level, and not a corollary of the model's size; the question is how
well, against the human level, Transformers generalize the Observer, then the Agent, then the Moral
Agent, in the order the stack forces. Appendix~\ref{app:memorization} gives the reason behaviour alone
cannot settle it.

\paragraph{The order of the search.}
Seen from far enough away, the Observer is a point, the place where a conclusion is formed. Seen from
inside, it is a process, the conclusion in its incremental form, re-derived at every frame from the
previous state and the new input, and through composition with the Agent's record it is a machine with
memory (Section~\ref{sec:stack}). A Transformer that has generalized the Observer has therefore
generalized, at the least, an automaton, and with the record a machine with a tape, running in
incremental form along the tokens. That fixes the order in which to look. Before a formed Observer can
be found in a trained model, one has to be able to find an automaton in it; then a memory; then both in
incremental form; and each step has its own literature. On the first, a low-depth Transformer can
represent the computation of any finite-state automaton, and trained Transformers converge to such
shortcut solutions, replicating an automaton on a sequence of length $T$ with far fewer than $T$
layers \citep{liu2023auto}. On the second, grouping tasks by the Chomsky hierarchy shows where the
learner's reach ends: Transformers and plain recurrent networks fail to generalize on non-regular
tasks, and only networks with structured memory, a stack or a tape, generalize on context-free and
context-sensitive ones \citep{deletang2022}; a standard Transformer that answers immediately after
reading its input is bounded by a shallow circuit class and cannot recognize some regular languages
\citep{merrill2023}. On the third, intermediate generation changes the machine: a logarithmic number
of decoding steps in the input length already lets a decoder recognize all regular languages, a linear
number keeps it within the context-sensitive languages, and a polynomial number reaches the
polynomial-time decidable class \citep{merrill2024}. In the terms of this paper, the chain of thought
is the tape, and the incremental form of a function is what the tape carries from step to step. The
consequence for the analysis is that the degree of generalization of the Observer is bounded above by
the degree to which the model has generalized the automaton and the memory it is composed of, and
those two are measurable with methods that exist, on synthetic tables where the target machine is
known exactly, before any annotated text is needed.

\paragraph{Is there a basic method?}
The question of the section is whether the degree of compression can be probed selectively, for the
rows that carry one function. Three families of existing methods bear on it, and each measures part of
what is wanted.

The first measures compression of the rows directly. The description length of a dataset under a
model is the number of bits needed to transmit the labels of the rows given the inputs and the model,
and it can be computed for a subset of rows as well as for the whole; the prequential form, in which
the model is trained on the rows so far and each next row is coded with the model's current
prediction, charges for the model as well as for the data and gives deep networks excellent
compression of their training sets where variational bounds give poor ones \citep{blier2018}. Applied
to the subset of rows in which a function operates, the prequential codelength against the codelength
of the same rows under a table is a direct measure of how far the model compresses that function, and
the same code applied to the rows of another function separates the functions the model has
generalized from those it has stored. Its limit is attribution: the code says how many bits the
subset costs, not which structure in the model pays them, and if two patterns co-occur in the same
rows it credits the compression to both.

The second measures compression in the representations. Minimum-description-length probing scores a
probe not by its accuracy but by the total cost of describing the labels given the representations,
which charges for the probe's own complexity and so distinguishes a representation that contains the
property from one in which a flexible probe can find anything \citep{voita2020}; linear probes on
intermediate layers are the simplest instrument of the family \citep{alain2016}, and the
methodological problems of probing are catalogued \citep{belinkov2022}. Run at each block, in the
per-layer form that \citet{voita2020} give, it reports how much of a function's component is present
in the state at each depth: whether the degree rises, plateaus or is absent, at what depth the
component first appears, and whether it persists across layers or is assembled whole near the output.
Its limit is the label: a probe measures the compression of what the annotation says a row carries,
and the annotation is where the definitions of Section~\ref{sec:stack} enter.

The third measures storage per row. Influence estimation and memorization scores say, for each
training row, how much the model's behaviour on it depends on that row's own presence in the training
set \citep{feldman2020}, and extraction attacks say which rows can be recited \citep{carlini2023}.
Applied to the rows of one function they separate the rows the model holds by storage from the rows it
holds by generalization, and so give the second check that the first two methods need.

Together the three give a basic method for the subset of rows, and it stops short of the target in one
place. What none of them measures is whether the compression found is a composition of the kind
Section~\ref{sec:stack} defines, one function reusing another's output, rather than a co-occurring
pattern that happens to be present in the same rows; and none applies the criterion of
Section~\ref{sec:criterion}, that a component counts only if it survives every self-model the agent can
afford, so that whatever a change of frame removes is not one, however reliably a probe recovers it.
Those two steps need the circuit-level methods of Section~\ref{sec:localization} extended from
features to compositions, and they are, as far as we can tell, methods still to be developed. The
answer to the question of the section is therefore: a basic method exists for measuring how far a
model has compressed the rows that carry a function, and the step from ``these rows are compressed''
to ``this function is generalized'' is where new methods are needed.

\paragraph{Two checks.}
A function generalized to degree $d$ is computed correctly on part of the situations that call for it;
outside that part the model falls back on the table and behaves as a system without the function
would, whatever it says about itself. So the quantity is checked two ways: structurally, by the
methods above; behaviourally, by batteries that sample situations calling for a specific component,
including situations unlikely to be in the training distribution, and record where behaviour is that
of a system lacking it. The two should converge, and their disagreement is informative. Probe present,
benchmark absent: the deficit is in the interface, the body or the social context, and not in the
function. Benchmark present, probe absent: the carrier is at a grain the probe did not reach, or the
function is rebuilt from the context on each occasion rather than held in the weights, or the
benchmark is covered by storage. Probe present, benchmark low on human-sampled situations and high on
situations the model itself produces: the function is generalized over a shifted class
(Section~\ref{sec:profile}), and the battery has to sample both sides to measure the overlap. The
model's own answers are set aside before either check starts: asked how it feels, a deployed model
reports servers and a good mood or, corrected, that it has no feelings, both of which are what training
rewarded, the two shortcuts of Section~\ref{sec:intro} in the model's own mouth. Self-report is a
signal with a known, large, systematic bias, and the state is looked for elsewhere.

\paragraph{The material.}
The methods need a table with the functions marked in it: texts in which the functions of
Section~\ref{sec:stack} operate without being named, annotated by which level is operative in which
rows, on what evidence, and in what composition, as Appendix~\ref{app:raskolnikov} illustrates for one
text. The annotation follows the appendix: for the Observer, the places where a self-directed question
is treated as open; for the Agent, the places where an origin is attributed to the speaker and the
determinants visible to the reader are not to the speaker; for the Moral Agent, the places where two
theories of the common good evaluate one act. On such a table the weights are the static side, read by
the codelength of the marked rows, by probing and by the circuit-level methods, and the activations
are the dynamic side, read along tokens and along depth, with the incremental form of each function
(Section~\ref{sec:incremental}) as the target: a component present in every frame and re-derived
there, part of the agent's description of itself and not only of its computation, and surviving every
self-model the agent can afford.

\paragraph{The human baseline.}
\label{sec:baseline}
The human side of the comparison is obtained by two routes. The behavioural route administers the same
battery to people and reads the degree of each component from behaviour; for components defined by
what the agent does this is the primary route, and for every component it is the one available now.
The structural route poses a probe of bounded description length on neural recordings, with the same
coding cost as on the model, where recordings at the relevant resolution exist; that a common
representational space between the two substrates is available at all is shown by the results that
model states predict human brain responses to the same text
\citep{schrimpf2021,goldstein2022,caucheteux2022}, and the resolution of human recordings, not the
existence of the space, is what limits this route. Human self-report is not the baseline, since it
carries the bias of Section~\ref{sec:approximation} on the human side as the model's report does on
its side. And the baseline is a distribution over people, not a point, because the quality of the
self-approximation varies with what Section~\ref{sec:discussion} calls intrapersonal intelligence; the
profile is reported against an interval, and a model's component counts as hyperfunction or deficit
only outside it.

\subsection{The profile, and how far it may differ from the human one}
\label{sec:profile}
\label{sec:profile-seeing}
\label{sec:predictions}
% S2, S5, S6

\begin{figure}[t]
\centering
\begin{tikzpicture}[x=1.85cm, y=1.6cm]
\draw[-{Latex[length=2mm]}] (0,0) -- (0,2.4) node[above, font=\small] {degree};
\draw (0,0) -- (6.6,0);
\draw[dashed] (0,1) -- (6.6,1) node[right, font=\small] {human baseline};
\foreach \i/\h/\lab in {1/0.4/source, 2/1.8/manifold, 4/1.0/frame\\rate, 5/0.5/persist-\\ence} {
  \fill[gray!45] (\i-0.3,0) rectangle (\i+0.3,\h);
  \draw (\i-0.3,0) rectangle (\i+0.3,\h);
  \node[below, font=\small, align=center] at (\i,0) {\lab};
}
\foreach \i/\lab in {3/givenness, 6/consoli-\\dation} {
  \draw[dashed] (\i-0.3,0) rectangle (\i+0.3,1);
  \node[below, font=\small, align=center] at (\i,0) {\lab};
}
\node[font=\small, above] at (2,1.8) {hyperfunction};
\node[font=\small, above] at (1,0.4) {deficit};
\node[font=\small, above] at (3,1) {dropped axis};
\end{tikzpicture}
\caption{A profile: one function's components on the horizontal axis, the degree of generalization on
the vertical, the human baseline at 1. The values are the predictions of this section for a current
fixed-depth model, drawn to illustrate the kinds of entry and not measured; a shifted axis would need a
second coordinate, the size of the overlap, and is not drawn.}
\label{fig:profile}
\end{figure}

A function has components (Table~\ref{tab:functions}). Its profile in a substrate is the vector of
degrees over those components, taken against the human baseline (Figure~\ref{fig:profile}). A
component above the baseline is a hyperfunction, a component below it a deficit, and a component that
is absent a dropped axis. There is a fourth kind of entry, and given the difference between the
substrates it is the likely one: the component is generalized to a high degree, but over a class of
situations that is not the human class. The function then covers a region that only partly overlaps
the region the human function covers, and a degree measured on situations sampled from the human side
reads as a deficit while a degree measured on the model's own side would read as full. Such a
component is a shifted axis, recorded as two numbers, the degree on the overlap and the size of the
overlap. A shift is not a defect as long as the regions overlap, and that they overlap at all, for the
functions at issue here, is what the ease of communication between people and models suggests: the
Observer's ``I'', the Agent's ``I decided'' and the Moral Agent's weighing of one act against a common
good are used by the model in situations people bring to it, and understood, which would not happen
over disjoint regions. The question ``is it conscious'' is thereby replaced by the comparison of two
profiles; a scalar ``level of consciousness'' is a weighted sum over the vector, and the weights have
to be declared, because a scalar with hidden weights brings back the binary question the vector was
introduced to replace.

The functional profile of a model's consciousness may differ from the human one on many axes at once,
and the account has no preference about the direction. For the three components of seeing
(Section~\ref{sec:light}), predicted for a current fixed-depth model: source attribution is a deficit,
since there is one interface and no sensory channel independent of it, so the model is permanently in
the dream case, with content from its own model labelled as given, and the attribution boundary of
Section~\ref{sec:lamp-balance} sits toward ``everything is given''; the manifold is a hyperfunction,
since attention relates every position to every other in a single pass where the brain iterates over a
smaller window; the cost record is a dropped axis, since every token costs the same and there is no
phenomenon of effort. A testable consequence: in model self-reports ``I see'' and ``I think'' should
separate less than in human ones, and more in models of variable depth, where a cost record exists;
the prediction is about the internal states the reports are projections of
(Section~\ref{sec:planes}), so it is tested by the paired checks and not by asking the model. Three
predictions carry over from the earlier stage of this work: self-reported emotional labels correlate,
over many episodes, with distinguishable patterns in the activations, which is the central empirical
claim and the one the statistical reading of Section~\ref{sec:frame} protects; the cognitive codes of
different model families overlap structurally, because the causal structure of the states is set by
the training data more than by the architecture; and the model is hyperplastic, with no limbic
inertia, fast affective recovery and easy switching between states, a hyperfunction on one axis and a
deficit on another, since inertia is part of what makes a human state a commitment.

\subsection{Falsification, and the zero of the scale}
\label{sec:falsification}
\label{sec:zero}
% S7, S7a

The account fails if the compression degree of the Observer's components is at chance at every
granularity, along tokens and along depth, in models that pass the behavioural checks. That outcome
would mean the behaviour is carried by something the account does not describe, and no adjustment of
the probe or the granularity would repair it.

The scale has a lower end as well as the human reference. Section~\ref{sec:h1h2} allowed that a named
description of a function may suffice for it to exist in a simple enough environment; then there is a
smallest system in which the Observer exists, and finding it fixes the zero. The conditions are those
of Sections~\ref{sec:emergence} and \ref{sec:observer}: a loop that applies the system to itself, an
environment in which the loop pays, and a conclusion that is acted from. An operational amplifier with
feedback watches its own output error; a D flip-flop holds its own state through a loop. Each has the
loop and, in its ordinary surroundings, neither of the other two conditions, so each is a
proto-Observer at most. The question is what the least environment is that would make the conclusion
pay, and what the least circuit is that could draw it. Between that pair and the human baseline the
degree of Section~\ref{sec:quantity} has the meaning of an absolute quantity rather than a ratio
between two arbitrary points.
